% Part: computability
% Chapter: machines-computations
% Section: configuration

\documentclass[../../../include/open-logic-section]{subfiles}

\begin{document}

\olfileid{cmp}{tur}{con}
\olsection{কনফিগারেশন ও গণনা}

\begin{explain}
আগে জোড়-যন্ত্রের কনফিগারেশনগুলি ধাপে ধাপে অনুসরণ করার কথা মনে করো।
ফিতা, পাঠ-লেখ হেড ও টুরিং-যন্ত্রের প্রোগ্রাম নিয়ে গঠিত কাল্পনিক
ব্যবস্থাটি আসলে টুরিং-যন্ত্রের গণনাকে চোখের সামনে দেখার একটি স্বজ্ঞাত
উপায়। আনুষ্ঠানিকভাবে, কোনো প্রদত্ত নিবেশে টুরিং যন্ত্রের গণনাকে
\emph{কনফিগারেশন}-এর একটি অনুক্রম হিসেবে সংজ্ঞায়িত করা যায়; আর একটি
কনফিগারেশন নিজে হলো প্রতীকের একটি অনুক্রম (গণনার কোনো এক মুহূর্তে ফিতার
বিষয়বস্তুর সঙ্গে যার মিল আছে), পাঠ-লেখ হেডের অবস্থান নির্দেশকারী একটি
সংখ্যা এবং একটি দশা। এগুলি ব্যবহার করে কোনো প্রদত্ত নিবেশে টুরিং
যন্ত্র~$M$ কী গণনা করে, তা সংজ্ঞায়িত করা যায়।
\end{explain}

\begin{defn}[কনফিগারেশন]
টুরিং যন্ত্র $M = \tuple{Q, \Sigma, q_0, \delta}$-এর একটি
\emph{কনফিগারেশন} হলো একটি ত্রয়ী $\tuple{C, m, q}$, যেখানে
\begin{enumerate}
\item $C \in \Sigma^*$ হলো $\Sigma$-র প্রতীকগুলির একটি সসীম অনুক্রম,
\item $m \in \Nat$ হলো একটি সংখ্যা $< \len{C}$, এবং
\item $q \in Q$।
\end{enumerate}
স্বজ্ঞাতভাবে, অনুক্রম~$C$ হলো ফিতার বিষয়বস্তু (বাঁদিকের প্রথম ঘর থেকে
শেষ যে ঘরটি ফাঁকা নয় অথবা আগে পরিদর্শিত হয়েছে, সেই ঘর পর্যন্ত সব
প্রতীক); পাঠ-লেখ হেডটি যে ঘর পড়ছে তার সংখ্যা হলো~$m$ (বাঁদিকের প্রথম
ঘরের সংখ্যা~$0$ থেকে শুরু); এবং $q$ হলো যন্ত্রটির বর্তমান দশা।
\end{defn}

\begin{explain}
টুরিং যন্ত্রের সম্ভাব্য নিবেশ হলো প্রতীকের একটি অনুক্রম—সাধারণত কোনোভাবে
একটি সংখ্যাকে সংকেতায়িত করা অনুক্রম। টুরিং যন্ত্রের আরম্ভিক কনফিগারেশন
হলো সেই কনফিগারেশন, যাতে ওই নিবেশের উপর কাজ করার জন্য যন্ত্রটিকে চালু
করি: ফিতায় শেষ-চিহ্নটির ঠিক পরে ডান দিকের ঘরগুলিতে নিবেশটি লেখা থাকে;
পাঠ-লেখ হেডটি নিবেশের জন্য নির্ধারিত বাঁদিকের প্রথম ঘরটি পড়ে (অর্থাৎ
বাঁ প্রান্তের চিহ্নটির ডান পাশের ঘর; শূন্য প্রতীকক্রমের ক্ষেত্রে ঘরটি
ফাঁকা); এবং ব্যবস্থাটি নির্দিষ্ট আরম্ভিক দশা~$q_0$-তে থাকে।
\end{explain}

\begin{defn}[আরম্ভিক কনফিগারেশন]
অশূন্য নিবেশ $I \in \Sigma^*$-এর জন্য $M$-এর \emph{আরম্ভিক
কনফিগারেশন} হলো
\[
\tuple{\TMendtape \frown I, 1, q_0}.
\]
নিবেশটি শূন্য প্রতীকক্রম হলে আরম্ভিক কনফিগারেশন হলো
\[
\tuple{\TMendtape \frown \TMblank, 1, q_0}.
\]
% BN-SRC-166 TARGET-CORRECTION-BEGIN
\par\smallskip\noindent\textnormal{\small\textbf{উৎস-সংশোধন (BN-SRC-166):} হিমায়িত ইংরেজি সংজ্ঞা সব $I\in\Sigma^*$-এর জন্য $\tuple{\TMendtape\frown I,1,q_0}$ দেয়। কিন্তু $I$ শূন্য প্রতীকক্রম হলে $C$-এর দৈর্ঘ্য~$1$, ফলে কনফিগারেশনের আগের শর্ত $m<\len{C}$ লঙ্ঘিত হয়। বাংলা পাঠে সেই ক্ষেত্রে প্রথম নিবেশ-ঘর হিসেবে একটি~$\TMblank$ যুক্ত করা হয়েছে।} % BN-SRC-166 TARGET-CORRECTION-END
\end{defn}

\begin{explain}
$\frown$ চিহ্নটি \emph{সংযুক্তি} বোঝায়—নিবেশ-প্রতীকক্রমটি বাঁ প্রান্তের
চিহ্নটির ঠিক ডান পাশে শুরু হয়।
% BN-SRC-165 TARGET-CORRECTION-BEGIN
\par\smallskip\noindent\textnormal{\small\textbf{উৎস-সংশোধন (BN-SRC-165):} হিমায়িত ইংরেজি ব্যাখ্যায় নিবেশটি বাঁ প্রান্তের চিহ্নটির ঠিক “বাঁ দিকে” শুরু হয় বলা হয়েছে। একমুখী ফিতা, আগের বর্ণনা এবং আরম্ভিক কনফিগারেশন—সব অনুযায়ী সেটি চিহ্নটির ঠিক ডান দিকে শুরু হয়; বাংলা পাঠে সেই দিকটি রাখা হয়েছে।} % BN-SRC-165 TARGET-CORRECTION-END
\end{explain}

\begin{defn}
আমরা বলি, একটি কনফিগারেশন $\tuple{C, m, q}$, $M$ অনুসারে
\emph{এক ধাপে $\tuple{C', m', q'}$ কনফিগারেশনটি দেয়}, যদি এবং কেবল যদি
\begin{enumerate}
\item $C$-এর $m$-তম প্রতীকটি~$\sigma$ হয়,
\item $M$-এর নির্দেশ-সেট $\delta(q, \sigma) =
  \tuple{q', \sigma', D}$ নির্দিষ্ট করে,
\item $C'$-এর $m$-তম প্রতীকটি~$\sigma'$ হয়, এবং
\item
\begin{enumerate}
\item $D = L$ এবং $m>0$ হলে $m' = m - 1$, নইলে $m'=0$; অথবা
\item $D = R$ এবং $m' = m + 1$; অথবা
\item $D = N$ এবং $m' = m$;
\end{enumerate}
\item $m' = \len{C}$ হলে $\len{C'} = \len{C} + 1$ এবং $C'$-এর
  $m'$-তম প্রতীকটি~$\TMblank$ হয়। অন্যথায় $\len{C'}=\len{C}$।
\item $i < \len{C}$ ও $i \neq m$ এমন সব~$i$-এর জন্য $C'(i) = C(i)$।
\end{enumerate}
\end{defn}

\begin{defn}\ollabel{defn:run-output}
নিবেশ~$I$-এ \emph{$M$-এর চালন} হলো $M$-এর কনফিগারেশনগুলির একটি সসীম
বা অসীম অনুক্রম~$C_i$, যেখানে $C_0$ হলো নিবেশ~$I$-এর জন্য $M$-এর
আরম্ভিক কনফিগারেশন, এবং অনুক্রমে থাকা প্রতিটি $C_{i+1}$-কে তার আগের
$C_i$ এক ধাপে দেয়।
% BN-SRC-167 TARGET-CORRECTION-BEGIN
\par\smallskip\noindent\textnormal{\small\textbf{উৎস-সংশোধন (BN-SRC-167):} হিমায়িত ইংরেজি সংজ্ঞা চালনকে শুধু “একটি অনুক্রম~$C_i$” বলে এবং প্রতিটি $C_i$-র পরে $C_{i+1}$ আছে ধরে নেয়। তাতে থামা চালনের কোনো শেষ কনফিগারেশন থাকতে পারে না। বাংলা পাঠে সসীম বা অসীম অনুক্রম এবং উপস্থিত পরপর দুই পদের শর্ত স্পষ্ট করা হয়েছে।} % BN-SRC-167 TARGET-CORRECTION-END

আমরা বলি, $M$ \emph{নিবেশ $I$-এ $k$ ধাপ পরে থামে}, যদি চালনটি
$C_k = \tuple{C, m, q}$-তে শেষ হয়, $C$-এর $m$-তম প্রতীকটি~$\sigma$
হয় এবং $\delta(q, \sigma)$ অসংজ্ঞায়িত হয়। সেক্ষেত্রে যদি
$\TMblank$-এ শেষ না-হওয়া এমন একটি প্রতীকক্রম~$O$ থাকে, যাতে কোনো
$j \in \Nat$-এর জন্য
$C = \TMendtape \concat O \concat \TMblank^j$, তবে নিবেশ~$I$-এর জন্য
$M$-এর \emph{নির্গম} হলো~$O$। ($\TMblank^j$ হলো $j$টি
$\TMblank$-এর একটি অনুক্রম।) এমন প্রতীকক্রম না থাকলে—যেমন শেষ-চিহ্নটি মুছে
ফেলা হয়ে থাকলে—এই নিয়মে থেমে যাওয়া চালনটির কোনো নির্গম নির্ধারিত হয়
না।
% BN-SRC-168 TARGET-CORRECTION-BEGIN
\par\smallskip\noindent\textnormal{\small\textbf{উৎস-সংশোধন (BN-SRC-168):} আগের অংশে হিমায়িত ইংরেজি পাঠ স্পষ্টভাবে $\TMendtape$-এর উপর অন্য প্রতীক লেখার অনুমতি দেয়, কিন্তু এখানে প্রতিটি থামা কনফিগারেশনকে নিঃশর্তভাবে $\TMendtape\concat O\concat\TMblank^j$ আকারের বলে ধরে নির্গমের অস্তিত্ব দাবি করে। বাংলা পাঠে ওই আকারটি থাকলেই নির্গম সংজ্ঞায়িত বলা হয়েছে এবং আকারটি না থাকলে পরিণতি স্পষ্ট করা হয়েছে।} % BN-SRC-168 TARGET-CORRECTION-END
\end{defn}

\begin{explain}
এই সংজ্ঞা অনুসারে, $M$-এর নির্গম~$O$ সবসময় $\TMblank$ ছাড়া অন্য কোনো
প্রতীকে শেষ হয়; আর সময়~$k$-তে বাঁদিকের প্রথম~$\TMendtape$ ছাড়া গোটা
ফিতা $\TMblank$ দিয়ে ভরা থাকলে $O$ হলো শূন্য প্রতীকক্রম।
\end{explain}

\end{document}
